% !PRINTED 364
\weberpage{364}

In order to derive from this the relation between $h$ and $h'$, we make use of the general group theorems of Bd.~II, \S\,2, which we formulate and supplement as follows:

12. \gesperrt{Theorem. If $A$ is a group and $B$ a divisor of $A$ of finite index $(A,B)$, and further $C$ a divisor of $B$ of finite index $(B,C)$, then, by Bd.~II, \S\,2, (4):}
\weq{4}{(A,C) = (A,B)(B,C).}
If $A$ is again a group with the divisor $B$ of index $\mu$, we set
\weq{5}{A = a_1 B + a_2 B + \cdots + a_\mu B,}
where $a_1, a_2, \dots, a_\mu$ is a complete system of representatives of $A$ with respect to $B$.

Now let $C$ be a group of elements which are composable with the elements of $A$ (e.\,g. in our case simply by multiplication); then from (5) there follows:
\weq{6}{AC = a_1 BC + a_2 BC + \cdots + a_u BC.} % [sic] German source prints u instead of \mu.
Now one and the same element can occur in the two subsidiary groups $a_1BC$, $a_2BC$ only if $a_1a_2^{-1}$ is contained in $BC$, but not in $B$.

From this we formulate the theorem:

13. \gesperrt{Theorem. If $A$ is a group, $B$ a divisor of $A$ of finite index $(A,B)$, and further $C$ a group composable with $A$ of such a kind that $A$ and $BC$ contain no common elements apart from those of $B$, then}
\weq{7}{(AC,\ BC) = (A,B).}
The assumption of this theorem can also be expressed thus:

\begin{center}\gesperrt{$B$ is the intersection of $A$ and $BC$,}\end{center}

and it is satisfied, e.\,g., when $C$ is contained in $B$, that is, when $B = BC$, and also when $C$ has no element in common with $A$.

14. \gesperrt{Theorem. If $A$ is a group composed of the two groups $A'$ and $B$}
\weq{8}{A = A'B}
\gesperrt{and $B'$ the intersection of $A'$ and $B$, then}
\weq{9}{(A,B) = (A',B'),}
\gesperrt{provided that these indices are finite.}

For by (8), in every subsidiary group $aB$ to $B$ in $A$ an element $a'$ from $A'$ is contained, and one can therefore represent these subsidiary
%CONT>>
